\chapter[Problem Formulation and Decision-Support Framework]%
        {Problem Formulation and\\Decision-Support \mbox{Framework}}
\label{ch:framework}

This chapter states the problem formally and describes the framework built to address it. It
sits between the literature reviewed in \Cref{ch:theory} and the methodological detail of
\Cref{ch:methodology}: the aim here is to make precise \emph{what} is being decided and
\emph{on what basis}, leaving \emph{how} each component works to the following chapter.

\section{The system under study}
\label{sec:system}

The operation modelled is a three-stage order fulfilment process. Orders arrive over the course
of a month, join a queue for picking, are picked, then queue for packing, are packed, then
queue for dispatch, and leave the system when dispatch is complete. \Cref{fig:flow} shows the
structure.

\begin{figure}[htbp]
\centering
\begin{tikzpicture}[node distance=0.42cm and 0.52cm]
  \node[note] (arr) {order\\arrivals};
  \node[queue, right=of arr] (q1) {$Q_1$};
  \node[stage, right=of q1]  (s1) {PICKING\\[-1pt]\footnotesize $w_1$ FTE};
  \node[queue, right=of s1]  (q2) {$Q_2$};
  \node[stage, right=of q2]  (s2) {PACKING\\[-1pt]\footnotesize $w_2$ FTE};
  \node[queue, right=of s2]  (q3) {$Q_3$};
  \node[stage, right=of q3]  (s3) {DISPATCH\\[-1pt]\footnotesize $w_3$ FTE};
  \node[note, right=of s3]   (done) {order\\complete};

  \foreach \a/\b in {arr/q1, q1/s1, s1/q2, q2/s2, s2/q3, q3/s3, s3/done}
    \draw[flow] (\a) -- (\b);

  % sequencing decision markers
  \foreach \q in {q1,q2,q3}
    \draw[dashed, draw=thesisblue!70] (\q.north) -- ++(0,0.55);
  \node[note, above=0.62cm of q2, text=thesisblue]
    {sequencing decision taken at every queue};

  \node[note, below=0.5cm of s2, text width=14cm]
    {unbounded queues $\cdot$ no blocking $\cdot$ workers dedicated to one stage $\cdot$
     non-pre-emptive service};
\end{tikzpicture}
\caption[Structure of the modelled fulfilment process]{Structure of the modelled fulfilment
process. Each stage has a dedicated worker pool and its own queue; the sequencing policy
governs which waiting order is taken next at each of the three queues.}
\label{fig:flow}
\end{figure}

Three features of this system generate the difficulty.

\textbf{Shared, finite capacity.} Each stage has a fixed number of workers for the month.
Workers are not transferable between stages, so capacity misallocated to packing cannot relieve
picking.

\textbf{Heterogeneous, stage-differentiated workload.} Orders differ in size and in product
characteristics, and those characteristics affect the three stages unequally. The same order
can be demanding at packing and undemanding at picking, or the reverse.

\textbf{Two service classes with different deadlines.} Urgent orders carry a much tighter
deadline than normal orders, and the two classes compete for the same workers at every stage.

\section{The decision problem}
\label{sec:decision-problem}

\subsection{Notation}

Let \(\mathcal{S} = \{1, 2, 3\}\) index the stages (picking, packing, dispatch) and
\(\mathcal{C} = \{u, n\}\) the urgency classes (urgent, normal). For a planning month:

\begin{center}
\begin{tabularx}{0.94\textwidth}{lX}
\toprule
\(\mathcal{O}\) & the set of orders to be processed in the month \\
\(u_{i,s}\) & workload units imposed by order \(i\) at stage \(s\) \\
\(t_{i,s}\) & realised service time of order \(i\) at stage \(s\) \\
\(a_i\) & arrival time of order \(i\) on the operating clock \\
\(c(i) \in \mathcal{C}\) & urgency class of order \(i\) \\
\(\sigma_c\) & system-time SLA threshold for class \(c\) \\
\(\mathbf{w} = (w_1, w_2, w_3)\) & workforce configuration, in monthly FTE \\
\(\pi \in \Pi\) & sequencing policy \\
\(H\) & operating horizon, in minutes \\
\(h, k\) & paid hours per worker-month, and hourly labour cost \\
\(\pi_c\) & penalty per late order of class \(c\) \\
\(\theta_c\) & required attainment floor for class \(c\) \\
\bottomrule
\end{tabularx}
\end{center}

An order \(i\) is on time if its system time \(f_i - a_i \leq \sigma_{c(i)}\), where \(f_i\) is
its dispatch completion time. Note that \(\sigma_c\) bounds the order's \emph{end-to-end time in
the system} --- arrival at picking to completion of dispatch, waiting included --- and is not a
bound on processing duration at any stage; \(t_{i,s}\) denotes the latter. Both are measured on
the operating clock of \Cref{sec:operating-time}. Orders not completed by \(H\) have undefined \(f_i\) and are
counted as failures. Class attainment is
\(\text{SLA}_c(\mathbf{w}, \pi) = \left|\{i : c(i) = c,\ \text{on time}\}\right| /
\left|\{i : c(i) = c\}\right|\).

\subsection{Statement of the problem}

The planning criterion is to choose a workforce configuration and a sequencing policy that
minimise expected total cost subject to both class-specific service floors being met:

\begin{equation}
\min_{\mathbf{w} \in \mathcal{W},\; \pi \in \Pi}\quad
C(\mathbf{w}, \pi)
\;=\;
\left(\textstyle\sum_{s} w_s\right) k h
\;+\; \sum_{c \in \mathcal{C}} \pi_c \, n^{\text{late}}_c(\mathbf{w}, \pi)
\label{eq:problem}
\end{equation}
\begin{equation}
\text{subject to}\quad
\text{SLA}_u(\mathbf{w}, \pi) \geq \theta_u,
\qquad
\text{SLA}_n(\mathbf{w}, \pi) \geq \theta_n.
\label{eq:problem-constraints}
\end{equation}

\Cref{eq:problem,eq:problem-constraints} state the \emph{planning criterion}: what a good plan
would be. They are not, and are not claimed to be, a programme the framework solves. The
essential difficulty is that \(\text{SLA}_c(\mathbf{w}, \pi)\) and
\(n^{\text{late}}_c(\mathbf{w}, \pi)\) have no closed form. They are emergent outcomes of a
stochastic queueing process with non-stationary arrivals, heterogeneous service requirements
and a sequencing decision taken at every stage. They can be evaluated, by simulation, but they
cannot be written down and optimised directly. The framework is therefore built around
\emph{evaluation of candidates} rather than solution of a program: \(\mathcal{W}\) is reduced
to a bounded, structured set \(\mathcal{W}_{\text{cand}}\) of operationally plausible
configurations, each of which is simulated under each policy.

What the implementation then does is apply an explicit lexicographic rule over the simulated
candidates. Writing \(\mathcal{F} = \{(\mathbf{w}, \pi) \in \mathcal{W}_{\text{cand}} \times \Pi
: \text{SLA}_u \geq \theta_u \wedge \text{SLA}_n \geq \theta_n\}\) for the subset that satisfies
\Cref{eq:problem-constraints}, and \(V(\mathbf{w},\pi) = \max(0, \theta_u - \text{SLA}_u) +
\max(0, \theta_n - \text{SLA}_n)\) for the violation score, the selected pair is

\begin{equation}
(\mathbf{w}^\star, \pi^\star) =
\begin{cases}
\arg\min_{(\mathbf{w},\pi)\,\in\,\mathcal{F}} C(\mathbf{w}, \pi)
  & \text{if } \mathcal{F} \neq \emptyset, \\[6pt]
\operatorname*{lexmin}_{(\mathbf{w},\pi)\,\in\,\mathcal{W}_{\text{cand}} \times \Pi}
  \bigl(V(\mathbf{w}, \pi),\, C(\mathbf{w}, \pi)\bigr)
  & \text{otherwise.}
\end{cases}
\label{eq:selection}
\end{equation}

The first branch enforces \Cref{eq:problem-constraints} exactly, restricted to the candidate
set. The second is a deliberate fallback: when no evaluated candidate meets both floors, the
framework does not refuse to answer --- it returns the candidate closest to feasibility, cheapest
among ties, and flags it as infeasible. That output is a diagnosis rather than a plan, and the
distinction is preserved wherever such a case is reported. \Cref{sec:selection-rule} states the
rule as implemented and \Cref{sec:adaptive} describes the local improvement step that follows it.

This is why the results of this thesis are never described as optimal. \Cref{eq:selection}
minimises over \(\mathcal{W}_{\text{cand}}\), a bounded structured neighbourhood, not over
\(\mathcal{W}\). They are the best outcomes found within an explicitly bounded search.

\subsection{Why the two decisions cannot be separated}
\label{sec:coupling}

It is worth being explicit about why \(\mathbf{w}\) and \(\pi\) must be chosen jointly rather
than in sequence, since the conventional practice is to fix one and then choose the other.

The coupling runs in both directions. In one direction, the workforce required to satisfy
\Cref{eq:problem-constraints} depends on the policy: a discipline that directs scarce capacity
at the orders whose deadlines bind will satisfy the urgent floor at a smaller \(\mathbf{w}\)
than one that ignores urgency entirely. Formally, the feasible set
\(\mathcal{F}(\pi) = \{\mathbf{w} : \text{SLA}_u \geq \theta_u \wedge \text{SLA}_n \geq
\theta_n\}\) is policy-dependent, and there is no reason to expect
\(\mathcal{F}(\pi_1) \subseteq \mathcal{F}(\pi_2)\) for any pair of policies. A workforce
sized under an assumed discipline may be infeasible under the discipline actually used.

In the other direction, the value of a policy depends on the workforce. A sequencing rule can
only act when there is a queue to choose from. As \(\mathbf{w}\) grows and utilisation falls,
decision points become rare, and in the limit every policy behaves identically because there is
usually only one order waiting. The difference \(C(\mathbf{w}, \pi_1) - C(\mathbf{w}, \pi_2)\)
is therefore not a constant to be estimated once; it is a function of \(\mathbf{w}\) that tends
to zero as capacity becomes ample. \Cref{sec:res-policy} measures exactly this.

A sequential procedure --- size the workforce from workload ratios, then pick a rule --- is
therefore solving neither problem correctly. It sizes capacity under an implicit policy
assumption it never states, and it then evaluates policies at a capacity that was chosen
without reference to them. The joint formulation of \Cref{eq:problem} avoids both errors at the
cost of requiring simulation to evaluate the objective.

\subsection{A worked example of the planning unit}
\label{sec:worked-example}

Because the semantics of the decision variable are the single most common source of
misinterpretation in capacity planning of this kind, it is worth working through one month
explicitly.

Take December. The month contains \num{40800} orders. Their expected picking workload, computed
deterministically from \Cref{eq:workload-units,eq:service-time}, is \num{204402} minutes of
picking work. One full-time equivalent supplies \(H = \SI{9600}{\minute}\) of productive
capacity per month. Sized at a target utilisation of \num{0.85}, the analytical estimate is

\begin{equation*}
\hat{w}_{\text{pick}} \;=\; \left\lceil \frac{\num{204402}}{\num{9600} \times 0.85} \right\rceil
\;=\; \left\lceil 25.05 \right\rceil \;=\; 26 \text{ FTE}.
\end{equation*}

The rounding here is itself instructive: the requirement exceeds 25 FTE by only two per cent of
one worker, and the ceiling operation converts that into a full additional hire. Applied
independently at three stages, this is the mechanism behind the persistent over-provision
reported in \Cref{sec:res-stage-gap}.

Simulation subsequently recommends 22. The interpretation of that number requires care on four
points.

First, \textbf{22 FTE is a quantity of labour, not a headcount on a shift.} It means 22
months of full-time picking capacity, or equivalently \(22 \times \num{9600} =
\num{211200}\) productive picking minutes available across the month. Whether that is delivered
by 22 full-time pickers, by 30 part-time pickers, or by 18 full-time pickers plus overtime, is
outside the model.

Second, \textbf{the \SI{160}{\hour} is productive time, not opening hours.} The warehouse is not
being modelled as open for \SI{160}{\hour}. The figure is the capacity one worker contributes in
a month, and the operation may well be open considerably longer.

Third, \textbf{the utilisation reported against that capacity is a ratio of busy minutes to
available minutes.} December's recommended configuration runs picking at
\SI{90.5}{\percent}, meaning that across the month, \SI{90.5}{\percent} of the purchased picking
capacity was spent in service rather than idle.

Fourth, \textbf{the service thresholds are measured on the same clock.} An urgent order's
\SI{240}{\minute} threshold is 240 minutes of the compressed operating timeline, not four hours
of wall-clock time from the customer's perspective.

Each of these follows directly from the identity between physical and paid capacity established
in \Cref{sec:design-principles}, and each is a statement the thesis is careful never to
contradict.

\subsection{What the formulation deliberately excludes}

Three simplifications in \Cref{eq:problem} should be stated, because each bounds the
interpretation of the result.

The workforce \(\mathbf{w}\) is an integer vector of monthly full-time equivalents, with no
intra-month structure. The formulation therefore cannot express the possibility that capacity
might be concentrated where demand is heaviest within the month. In an operation with a
pronounced weekly cycle this is a real restriction, and it means the recommendation should be
read as a monthly capacity envelope rather than a deployment plan.

Workers are dedicated to a stage. There is no term allowing capacity to be reallocated between
picking, packing and dispatch, either in advance or dynamically. Since the results show pressure
concentrating almost entirely at one stage, cross-training would plausibly be valuable, and its
absence means the framework may over-state the total workforce required.

The cost function is linear in workers and in late orders. Real labour costs are not linear ---
recruitment, training and overtime introduce thresholds and non-linearities --- and contractual
penalties are frequently structured rather than per-order. The linear form is a transparent
approximation whose parameters are stated explicitly and can be varied.

\section{Framework architecture}
\label{sec:architecture}

The framework is organised as a pipeline, shown in \Cref{fig:architecture}. Its two entry
points --- an order-level history or an aggregate forecast --- converge after order
construction, so that everything downstream is shared between the two workflows.

\begin{figure}[htbp]
\centering
\begin{tikzpicture}[node distance=0.52cm]
  \node[entry] (hist) {Retrospective\\[-1pt]\footnotesize order-level history};
  \node[entry, right=1.1cm of hist] (fut)
       {Forecast\\[-1pt]\footnotesize aggregate volume $+$ uncertainty};

  \node[block, below=0.85cm of $(hist)!0.5!(fut)$] (orders)
       {Orders carrying stage-specific workload units};
  \node[block, below=of orders] (compress)
       {Compression onto the operating horizon $H = \SI{9600}{\minute}$};
  \node[block, below=of compress] (cand)
       {Analytical capacity estimate $\rightarrow$ 16 workforce candidates};
  \node[block, below=of cand, fill=thesisblue!7, draw=thesisblue] (sim)
       {\textbf{Simulation:} each candidate $\times$ each policy,\\
        under common random numbers};
  \node[block, below=of sim] (feas)
       {Class-specific feasibility screen $+$ economic evaluation};
  \node[block, below=of feas] (diag)
       {Stage Pressure diagnostic $\rightarrow$ adaptive capacity test};
  \node[endnode, below=of diag] (rec)
       {Recommendation: workforce, policy, cost, service, bottleneck};

  \draw[flow] (hist.south) -- ++(0,-0.28) -| (orders.north);
  \draw[flow] (fut.south)  -- ++(0,-0.28) -| (orders.north);
  \foreach \a/\b in {orders/compress, compress/cand, cand/sim, sim/feas, feas/diag, diag/rec}
    \draw[flow] (\a) -- (\b);

  \node[note, right=0.35cm of sim, text width=2.6cm, anchor=west]
       {shared by\\both workflows};
  \begin{scope}[on background layer]
    \node[fit=(orders)(rec), draw=black!25, dashed, rounded corners,
          inner sep=5pt] {};
  \end{scope}
\end{tikzpicture}
\caption[Framework architecture]{Architecture of the decision-support framework. The two entry
points converge after order construction, so both workflows share the same simulation engine,
feasibility criterion, economic model and diagnostics (dashed region).}
\label{fig:architecture}
\end{figure}

\section{Design principles}
\label{sec:design-principles}

Five principles shaped the design, and each is defended in \Cref{ch:methodology}.

\textbf{Physical capacity equals paid capacity.} The hours a worker can process work in
simulation are exactly the hours the economic model pays for. Without this identity, capacity
recommendations are systematically optimistic.

\textbf{The horizon is finite and backlog is explicit.} A month is a month. Work not finished
within it is retained as visible backlog rather than absorbed by unlimited overtime.

\textbf{Feasibility is class-specific.} A configuration must satisfy both service floors
separately. Aggregate attainment can conceal complete failure of a minority class.

\textbf{Comparisons are made on identical stochastic realisations.} Policies are compared under
common random numbers, so observed differences reflect the sequencing decision rather than
sampling noise.

\textbf{Recommendations must be explainable.} Every recommendation is accompanied by the
identity of the most pressured stage, the decomposition of that judgement, and an explicit
economic test of whether relieving it would be worthwhile. A number a planner cannot interrogate
is not a decision support.

\section{The two workflows}
\label{sec:workflows}

\subsection{Retrospective analysis}

The retrospective workflow answers a counterfactual question: \emph{given the demand that
actually occurred, which workforce configuration and sequencing policy would this framework
have recommended?} It consumes an order-level dataset, processes each month independently, and
produces a recommendation per month.

The framing matters. The output is not a claim that the recommended workforce was the one
actually deployed, nor that deploying it would have produced exactly the reported outcome. It
is a statement about what the model, under its stated assumptions, would have advised. Its
value is diagnostic: it reveals how required capacity varies across a demand year, where
pressure concentrates, and how much the sequencing discipline matters at each point in the
seasonal cycle.

\subsection{Forecast-based planning}

The forecast workflow answers a prospective question: \emph{given an expected demand volume for
a coming month, what workforce should be planned?} It consumes an aggregate forecast rather
than individual orders.

The distinction is important and easily overstated in the wrong direction. The framework does
\emph{not} predict which orders will arrive, and it does not reuse the order-level dataset the
retrospective workflow consumes. It transforms an aggregate demand assumption into newly
generated, plausible simulated order streams that are consistent with that assumption and with
the configured seasonal, urgency and product-mix structure, then evaluates candidate workforces
against those streams. Because the streams are stochastic, the workflow generates several
replications and reports evidence across them --- a mean cost, a measure of the spread observed
and the share of replications in which the service floors were met --- rather than a single
deterministic answer. With the three replications used here, that evidence is indicative of
spread rather than a quantified estimate of risk, a limitation \Cref{sec:res-future} states
explicitly.

The two workflows therefore have genuinely different uncertainty structures. The retrospective
mode holds demand fixed and asks what would have been advisable; the forecast mode treats
demand as uncertain and asks what is robust. \Cref{sec:disc-modes} returns to what this means
for how their outputs should be used.

\subsection{What each workflow may and may not be claimed to do}

Because both workflows produce a recommendation in the same format, it is easy to over-read
either one. Two boundaries are worth stating at the outset.

The retrospective workflow performs \emph{counterfactual} analysis. Given demand that has
already occurred, it identifies the configuration the model would have advised. It makes no
claim that this configuration was the one deployed, that deploying it would have produced
exactly the reported outcome, or that the outcome observed in reality is comparable. Its proper
use is diagnostic --- to reveal how the capacity requirement moves across a demand year, where
pressure concentrates, and under which conditions the operating discipline matters --- rather
than evaluative of any past decision.

The forecast workflow does \emph{not} predict individual future orders. This distinction is
easy to blur and important to preserve. The workflow takes an aggregate assumption about
volume, and generates order streams consistent with that assumption and with the configured
structural profile of demand. The resulting orders are plausible instances, not predictions.
What the workflow therefore evaluates is whether a candidate configuration holds up across a
range of demand realisations that the assumption admits --- which is a statement about
robustness under an assumed volume, not a forecast of what will happen. If the volume
assumption itself is wrong, nothing in the workflow will detect it.
